\documentclass[11pt]{article}

\usepackage[utf8]{inputenc}
\usepackage[T1]{fontenc}
\usepackage{geometry}
\usepackage{amsmath}
\usepackage{amssymb}
\usepackage{booktabs}
\usepackage{hyperref}
\usepackage{xurl}
\usepackage{xcolor}
\usepackage{listings}
\usepackage{enumitem}
\usepackage{parskip}
\usepackage{graphicx}
\graphicspath{{../images/}}

\geometry{margin=1in}

\lstset{
  language=Java,
  basicstyle=\ttfamily\small,
  keywordstyle=\color{blue},
  commentstyle=\color{gray},
  stringstyle=\color{teal},
  showstringspaces=false,
  columns=fullflexible,
  frame=single,
  framerule=0.2pt,
  breaklines=true
}

\setlist[itemize]{topsep=0.3em,itemsep=0.2em}
\setlist[enumerate]{topsep=0.3em,itemsep=0.2em}

% Simple per-section source line.
\newcommand{\notesources}[1]{\par\noindent{\small\color{gray}\textit{Sources:} \sloppy #1\par}}

% Unit-wise source strings for this subject (used by slides/notes).
% Path is relative to the lecture `latex/` folder (the compilation CWD).
% OOPS with Java (CSEG1044) — unit-wise sources
%
% These are short, student-facing reference strings intended to be used with:
%   - \slidesources{...}  (in beamer slides)
%   - \notesources{...}   (in lecture notes)
%
% Style inspiration: other subjects in My-Subjects/ use semicolon-separated sources
% and include an "accessed" date for web documentation.

\newcommand{\OOPSUnitISources}{Schildt, \textit{Java: A Beginner's Guide} (10e), McGraw-Hill (Java fundamentals + OOP basics).; Oracle Java Tutorials: Getting Started + Language Basics + Classes and Objects (accessed \today).; Java SE API docs (e.g., \texttt{String}, \texttt{Scanner}) (accessed \today).}

\newcommand{\OOPSUnitIISources}{Schildt, \textit{Java: The Complete Reference} (12e), McGraw-Hill (classes, inheritance, packages, interfaces).; Bloch, \textit{Effective Java} (3e), Addison-Wesley, 2018 (design choices; interfaces vs abstract classes).; Oracle Java Tutorials: Classes and Objects + Interfaces + Packages (accessed \today).; Java SE API docs (e.g., \texttt{Comparable}, \texttt{Runnable}) (accessed \today).}

\newcommand{\OOPSUnitIIISources}{Schildt, \textit{Java: The Complete Reference} (12e) (nested/inner classes, exceptions, threads, I/O).; Oracle Java Tutorials: Nested Classes + Exceptions + Concurrency + Basic I/O (accessed \today).; Java SE API docs (e.g., \texttt{Thread}, \texttt{AutoCloseable}) (accessed \today).}

\newcommand{\OOPSUnitIVSources}{Schildt, \textit{Java: The Complete Reference} (12e) (generics, lambdas, Swing, JDBC).; Bloch, \textit{Effective Java} (3e) (generics + lambdas).; Oracle Java Tutorials: Generics + Lambda Expressions + Swing + JDBC Basics (accessed \today).; Java SE API docs (e.g., \texttt{java.util.function}, \texttt{javax.swing}, \texttt{java.sql}) (accessed \today).}

\newcommand{\OOPSUnitVSources}{Schildt, \textit{Java: The Complete Reference} (12e) (Collections Framework + wrapper classes).; Oracle Java docs: Collections Framework overview (accessed \today).; Java SE API docs (\texttt{java.util} collections; wrapper classes in \texttt{java.lang}) (accessed \today).}

\newcommand{\OOPSUnitVISources}{Git SCM: \textit{Pro Git} (2e) (version control workflow), accessed \today.; GitHub Docs (repositories, issues, pull requests), accessed \today.; Oracle Java Tutorials: Swing + JDBC (accessed \today).; JUnit 5 User Guide (accessed \today).; Apache Maven Project: Getting Started (accessed \today).; Gradle User Manual: Getting Started (accessed \today).; UML class diagrams: UML 2.x overview/spec (accessed \today).}

\newcommand{\OOPSMiscWhyInterfacesSources}{Schildt, \textit{Java: The Complete Reference} (12e) (interfaces).; Bloch, \textit{Effective Java} (3e) (prefer interfaces where appropriate).; Oracle Java Tutorials: Interfaces (accessed \today).; Java SE API docs (e.g., \texttt{Comparable}, \texttt{Runnable}) (accessed \today).}



\title{Object Oriented Programming (CSEG1044)\\Unit 02 -- Lecture 01 Notes\\Classes, Objects, and Encapsulation (Basics)}
\author{Tofik Ali}
\date{\today}

\begin{document}
\maketitle
\tableofcontents

\section{Lecture Overview}
In Unit 01 we wrote programs using variables, control flow, arrays, and strings.
Now we start writing \textbf{real OOP code}: we will create classes and objects and learn how encapsulation makes programs safer.
\notesources{\OOPSUnitIISources}

\section{Syllabus Alignment (Unit 02)}
\textbf{Unit 02:} Classes, Inheritance, Packages and Interfaces\par
\textbf{Today:} classes/objects + encapsulation basics

\section{Learning Outcomes}
After this lecture, you should be able to:
\begin{itemize}
  \item Identify fields vs methods and explain ``state vs behavior'' for a class.
  \item Explain object references in Java (aliasing) using a small example.
  \item Apply encapsulation using \texttt{private} fields and validation methods.
  \item State at least one invariant for a class and protect it in code.
\end{itemize}

\section{Key Terms (Quick)}
\begin{itemize}
  \item \textbf{Field}: data stored inside an object (state).
  \item \textbf{Method}: behavior (operations) of an object.
  \item \textbf{Reference}: a variable that points to an object (not the object itself).
  \item \textbf{Aliasing}: two references pointing to the same object.
  \item \textbf{Encapsulation}: keeping fields private and changing state through methods.
  \item \textbf{Access modifier}: \texttt{private}, default, \texttt{protected}, \texttt{public} (controls visibility).
\end{itemize}

\section{Detailed Notes}
\subsection{Class fundamentals: fields, methods, state vs behavior}
A \textbf{class} is a blueprint. It defines:
\begin{itemize}
  \item \textbf{Fields} (what data is stored),
  \item \textbf{Methods} (what operations are allowed),
  \item and usually a \textbf{constructor} (how objects are created).
\end{itemize}

\textbf{State vs behavior:}
\begin{itemize}
  \item State = current values in fields (e.g., \texttt{cgpa=8.2}).
  \item Behavior = operations that read/change state (e.g., \texttt{updateCgpa(...)}).
\end{itemize}

\paragraph{Rule of thumb.} A class should have a clear responsibility. If it tries to do everything, it becomes hard to maintain.

\subsection{Objects and reference semantics (mental model)}
In Java, objects are created on the heap (conceptually), and variables hold \textbf{references} to objects.
So assignment copies the reference (not the whole object).

\textbf{Aliasing example:}
\begin{itemize}
  \item \texttt{Student a = new Student(...);}
  \item \texttt{Student b = a;} $\rightarrow$ now both \texttt{a} and \texttt{b} refer to the same object.
  \item Changes through \texttt{b} are visible through \texttt{a}.
\end{itemize}

\paragraph{Why it matters.} Many bugs come from unexpected aliasing. Later, you will use immutability and defensive copying in some designs.

\subsection{Encapsulation: protect invariants}
Encapsulation is not only ``hiding data''; it is mainly about \textbf{protecting invariants}.
If fields are public, any code can set invalid values.

\textbf{Example invariant:} CGPA must be in range \texttt{0.0} to \texttt{10.0}.

\textbf{Approach:}
\begin{itemize}
  \item Make fields \texttt{private}.
  \item Provide methods that validate inputs before updating the state.
  \item Expose only what users need (avoid unnecessary setters).
\end{itemize}
\notesources{\OOPSUnitIISources}

\section{Worked Example: \texttt{Student} (encapsulation + references)}
\subsection*{Code}
\begin{lstlisting}
public class Student {
    private final String rollNo;
    private final String name;
    private double cgpa; // invariant: 0.0 <= cgpa <= 10.0

    public Student(String rollNo, String name, double cgpa) {
        if (rollNo == null || rollNo.isBlank()) throw new IllegalArgumentException("rollNo");
        if (name == null || name.isBlank()) throw new IllegalArgumentException("name");
        if (cgpa < 0.0 || cgpa > 10.0) throw new IllegalArgumentException("cgpa");
        this.rollNo = rollNo;
        this.name = name;
        this.cgpa = cgpa;
    }

    public double getCgpa() { return cgpa; }

    public void updateCgpa(double newCgpa) {
        if (newCgpa < 0.0 || newCgpa > 10.0) throw new IllegalArgumentException("cgpa");
        this.cgpa = newCgpa;
    }

    public String summary() {
        return rollNo + \" - \" + name + \" (CGPA=\" + cgpa + \")\";
    }
}
\end{lstlisting}

\subsection*{Aliasing demo (main)}
\begin{lstlisting}
public class Demo {
    public static void main(String[] args) {
        Student a = new Student("21CS1001", "Asha", 8.2);
        Student b = a;                 // aliasing
        b.updateCgpa(9.0);             // changes same object
        System.out.println(a.getCgpa()); // prints 9.0
    }
}
\end{lstlisting}

\section{In-Class Exercises (with brief solutions)}
\begin{enumerate}
  \item Two references \texttt{a} and \texttt{b} point to the same object. If \texttt{b} modifies it, does \texttt{a} see the change?\par
  \textbf{Answer:} yes (aliasing).

  \item Why is it risky to keep fields \texttt{public}?\par
  \textbf{Answer:} any code can put the object into an invalid state, breaking invariants.

  \item Write one invariant for a \texttt{BankAccount} class.\par
  \textbf{Sample:} \texttt{balance >= 0}.
\end{enumerate}

\section{Self-Study Practice (no solutions)}
\begin{enumerate}
  \item Design a \texttt{LibraryBook} class with an invariant (e.g., cannot issue if already issued).
  \item Write a small program that shows aliasing using a class of your choice.
  \item Refactor a ``public fields'' class into an encapsulated design with validation methods.
\end{enumerate}

\section{Checkpoint Questions (with sample answers)}
\textbf{Q1.} What is encapsulation and why do we use it?\par
\textbf{Sample answer:} Encapsulation means keeping internal data private and changing it only through methods. It helps protect invariants and prevents accidental misuse, making code easier to maintain.

\vspace{0.6em}
\textbf{Q2.} What is aliasing?\par
\textbf{Sample answer:} Aliasing happens when two references point to the same object. Updates through one reference are visible through the other.

\end{document}
